\documentclass[a4paper]{report}

% =========================================
% PAGE SETUP
% =========================================
\usepackage[a4paper]{geometry}
% =========================================
% BASIC PACKAGES
% =========================================
\usepackage{graphicx}       % Images
\usepackage{amsmath}        % Mathematics
\usepackage{amssymb}        % Mathematical symbols
\usepackage{booktabs}       % Professional tables
\usepackage{float}          % [H] figure/table positioning
\usepackage{caption}        % Caption customization
\usepackage{hyperref}       % Hyperlinks and references
\usepackage{ragged2e}       % Text alignment
\usepackage{setspace}       % Line spacing
\usepackage{hyperref}       % Clickable links and references
\usepackage{pifont}
\usepackage{tabularx}
\usepackage{array}
\usepackage[retainorgcmds]{IEEEtrantools}
\usepackage{algorithm}
\usepackage{algpseudocode}
% =========================================
% APA REFERENCING
% =========================================

\usepackage[
    style=apa,
    backend=biber
]{biblatex}

\addbibresource{references.bib}


% =========================================
% HYPERLINK SETTINGS
% =========================================
\hypersetup{
    colorlinks=true,
    linkcolor=black,
    urlcolor=black,
    citecolor=black,
    pdftitle={Critical Research Review Article}
}

% Table and figure caption formatting
\DeclareCaptionFont{captionnine}{\fontsize{9pt}{11pt}\selectfont}

\captionsetup[table]{
    font=captionnine,
    justification=raggedright,
    singlelinecheck=false,
    labelsep=colon
}

\captionsetup[figure]{
    font=captionnine,
    justification=centering,
    singlelinecheck=false,
    labelsep=colon
}


% =========================================
% DOCUMENT
% =========================================
\begin{document}


\begin{center}

{\fontsize{18}{22}\selectfont\textbf{RESEARCH REPORT}}

\vspace{1cm}

{\fontsize{18}{22}\selectfont\textbf{2D PROJECTILES SIMULATION\\[6pt]}}

\vspace{1cm}

{\fontsize{12}{14}\selectfont\textbf{Authors}}

\vspace{0.3cm}

{\fontsize{12}{14}\selectfont
Elsa Z. Fakhurrazi\\
}

\vspace{0.8cm}

{\fontsize{12}{14}\selectfont
\textbf{Project:} Documentation on the investigation of the fundametals laws of motion\\
\vspace{0.8cm}
\textbf{Project Version:} v1 \\
\textbf{Supervisor:} Father\\
\textbf{Start Date:} 7th October 2026\\
\textbf{Date of Completion:} N/A

}

\end{center}

% =========================================
% ABSTRACT
% =========================================

\vspace{0.5cm}

\noindent
\textbf{Abstract}

nothing yet 


\vspace{0.2cm}

\begin{justify}



\vspace{0.3cm}

\end{justify}



% =========================================
% TABLE OF CONTENTS
% =========================================

\newpage

\tableofcontents

% =========================================
% MAIN CONTENT
% =========================================


% =========================================
% 1. INTRODUCTION
% =========================================
\newpage

\section{Introduction}

The aim of this report is to document the process of my interpretation and derivation of the fundamnetal laws of motions. More specifically, 
the projectile motions of a simple object. I intend to document each step of my simulation and formula derivation to 
deepen my understanding of the mechanics of motions in the physical reality. 

\section{Fundamental Formulas Of Motions}

The following formulas is used to predict an object's acceleration, time, displacement and etc, by providing a combination of values from other variables. 

\begin{equation}
    v = u + at
\end{equation}
\begin{equation}
    s = ut + \frac{1}{2}at^2 
\end{equation}
\begin{equation}
    v^2 = u^2 + 2as
\end{equation}
\subsection{Derivation of the first formula}

Acceleration is defined as the rate of change in velocity. In other words, how fast the speed of the object changes. To measure this,
We take the difference between the final velocity with the initial velocity, divided by time, to see how much per second, does the speed change over time. 
As such, the formula is as follows:

\begin{equation}
    a = \frac{v - u}{t}
\end{equation}

The derivation of the first formula is simply a rearrangement of the acceleration formula, creating velocity as the subject.

\begin{IEEEeqnarray}{rCl}
    a  & = & \frac{v - u}{t} \\
    at & = & v - u \\
    v  & = & u + at
\end{IEEEeqnarray}



\end{document}

